\chapter{从综合得分到约束型多头组合}
\label{ch:portfolio}

\section{本章导读与学习目标}

第 \cref{ch:factor-combination} 已经输出唯一冻结的 \texttt{combination\_version}、证券层面综合得分、共同样本状态和适用条件。本章把这份横截面排序信息转化为信号日结束后“希望持有”的多头目标权重，并生成相对上一期实际持仓的期望交易差额。核心问题是：\keyterm{如何在不偷看下一交易日的前提下，让目标组合既表达综合分，又满足基础约束并可被完整审计？}

完成本章后，读者应能够：

\begin{enumerate}
  \item 区分综合得分、目标组合、期望订单、实际成交和实际持仓；
  \item 用冻结的固定数量规则和确定性并列规则构造候选集合；
  \item 从基础等权出发，依次处理个股、行业、流动性和换手约束；
  \item 解释持仓缓冲、现金权重与目标换手的作用和边界；
  \item 在约束冲突时执行预注册、确定性的降级流程；
  \item 输出可追溯的目标权重、期望交易差额、约束状态和组合版本；
  \item 为第 10 章保留目标与实际之间的差异，而不提前模拟成交。
\end{enumerate}

\begin{principlebox}
\textbf{[通用原则]} 下一交易日无法成交不能反向改变信号日综合分、候选名单或原始目标权重。未成交和实际持仓偏离应由第 10 章的订单、成交与持仓账本处理。
\end{principlebox}

\begin{figure}[htbp]
  \centering
  \begin{tikzpicture}[
    node distance=2.5mm,
    every node/.style={font=\footnotesize},
    quantPortfolioScore/.style={quantBlueCard, text width=2.15cm, minimum height=1.15cm, align=center},
    quantPortfolioCandidate/.style={quantCard, text width=2.15cm, minimum height=1.15cm, align=center},
    quantPortfolioBase/.style={quantCard, text width=2.15cm, minimum height=1.15cm, align=center},
    quantPortfolioConstraint/.style={quantOrangeCard, text width=2.15cm, minimum height=1.15cm, align=center},
    quantPortfolioTarget/.style={quantGreenCard, text width=2.15cm, minimum height=1.15cm, align=center},
    quantPortfolioTrade/.style={quantBlueCard, text width=2.15cm, minimum height=1.15cm, align=center}
  ]
    \node[quantPortfolioScore] (score) {冻结综合分};
    \node[quantPortfolioCandidate, right=of score] (candidate) {候选集合};
    \node[quantPortfolioBase, right=of candidate] (base) {基础等权};
    \node[quantPortfolioConstraint, right=of base] (constraint) {个股、行业、流动性、换手};
    \node[quantPortfolioTarget, right=of constraint] (target) {目标权重与现金};
    \node[quantPortfolioTrade, right=of target] (trade) {期望交易差额};

    \draw[quantArrow] (score) -- (candidate);
    \draw[quantArrow] (candidate) -- (base);
    \draw[quantArrow] (base) -- (constraint);
    \draw[quantArrow] (constraint) -- (target);
    \draw[quantArrow] (target) -- (trade);
  \end{tikzpicture}
  \caption{从冻结综合分到期望交易差额的规则型流程（结构示意）}
  \label{fig:ch09-portfolio-flow}
\end{figure}

本章不判断下一交易日真实成交，不处理停牌、价格限制、部分成交、正式交易成本、实际持仓更新或组合收益。这些全部留给第 10 章。

\section{为什么综合得分不是持仓}

\texttt{combined\_score} 只表达证券在冻结因子体系下的相对优先级，没有资金单位，也没有自动包含行业、流动性、换手和个股集中度。把分数直接归一化为权重，可能产生负权重、集中于少数证券，或让某一共同暴露被放大。

目标权重是一项带约束的决策。它既要尽可能保留排序信息，又要接受可行性边界。约束可能削弱原始信号表达；这种损失应通过基础权重、每一步调整和最终状态被记录，而不是隐藏在一个最终权重字段中。

\section{综合分、目标权重、订单和实际持仓的区别}

\begin{table}[htbp]
  \centering
  \small
  \caption{五类对象的职责与时间边界}
  \label{tab:ch09-object-responsibilities}
  \begin{tabularx}{\textwidth}{p{2.7cm} p{4.5cm} p{3.2cm} Y}
    \toprule
    对象 & 表达含义 & 形成时点 & 本章是否生成 \\
    \midrule
    综合得分 & 证券横截面排序信息 & 信号日数据冻结后 & 读取，不改写 \\
    目标组合 & 信号日后希望持有的权重 & 约束构建完成后 & 是 \\
    期望订单 & 目标与上一期实际持仓之差 & 目标权重确定后 & 是，仅目标差额 \\
    实际成交 & 下一交易日真实执行结果 & 交易状态与价格可得后 & 否 \\
    实际持仓 & 成交后真实持有状态 & 成交账本更新后 & 否 \\
    \bottomrule
  \end{tabularx}
\end{table}

期望交易为正表示希望增加，为负表示希望减少，为零表示目标不变，但都不代表已成交。若卖单在下一交易日无法成交，旧持仓仍保留在实际账本中；不能因目标权重为零就让它凭空消失。

\section{本章冻结输入和禁止访问的数据}

本章只读取第 8 章冻结的 \texttt{combination\_version} 及其 \texttt{signal\_date}、\texttt{security\_id}、\texttt{combined\_score}、\texttt{joint\_valid} 和 \texttt{applicability\_conditions}；同时读取 \cref{ch:universe-labels} 的信号日股票池、信号日点时行业、市值与历史流动性，以及上一期实际持仓或期初快照。

禁止访问的字段包括 \texttt{execution\_status}、\texttt{filled}、\texttt{execution\_price}、\texttt{limit\_status}、\texttt{next\_day\_suspension}、\texttt{future\_return} 和 \texttt{backtest\_result}。若输入表意外携带这些字段，目标构建应立即失败，而不是简单忽略后继续运行。

\begin{confirmbox}
\textbf{[需实际确认]} 真实资金规模、基准行业权重、行业偏离、参与率、个股上限、持股数量、换手预算和持仓快照截止时点必须由版本化配置确认。本章不虚构团队阈值。
\end{confirmbox}

\section{从综合分到候选集合}

在信号日 \(t\)，先保留 \texttt{joint\_valid = True} 且满足冻结适用条件的证券，再按 \texttt{combined\_score} 从高到低排序。教学主基线使用预注册的固定数量 \(N\)，候选集合可写为

\begin{equation}
  C_t
  = \operatorname{TopN}_{i\in U_t^{\mathrm{eligible}}}
    \left(\operatorname{combined\_score}_{i,t};N\right).
  \label{eq:ch09-candidate-set}
\end{equation}

并列时依次使用冻结得分、预注册的次级排序字段和稳定证券主键。稳定主键只用于确定性打破并列，不表达经济偏好。每条记录都保存 \texttt{score\_rank}、\texttt{selected} 与 \texttt{selection\_reason}，未入选记录仍保留。

\section{固定数量与固定分位入选规则}

\begin{table}[htbp]
  \centering
  \small
  \caption{固定数量与固定分位规则的比较}
  \label{tab:ch09-selection-rules}
  \begin{tabularx}{\textwidth}{p{2.6cm} Y Y}
    \toprule
    比较项 & 固定数量 & 固定分位 \\
    \midrule
    持股数量 & 较稳定，便于解释与审计 & 随有效股票池规模变化 \\
    信号池变化 & 入选比例随池规模变化 & 入选比例相对稳定 \\
    小样本风险 & 候选不足时直接返回状态 & 分位边界和取整更敏感 \\
    并列处理 & 第 \(N\) 名附近需确定性规则 & 分位阈值附近需确定性规则 \\
    本章角色 & 教学主基线 & 预注册敏感性方案 \\
    \bottomrule
  \end{tabularx}
\end{table}

持股数量必须在查看第 10 章正式回测前冻结。不能依据净收益、成本后表现或成交成功率在两种规则之间择优。

\section{基础等权组合}

设候选集合为 \(C_t\)，约束前基础权重为

\begin{equation}
  w^{\mathrm{base}}_{i,t}
  =
  \begin{cases}
    1/|C_t|, & i\in C_t,\\
    0, & i\notin C_t.
  \end{cases}
  \label{eq:ch09-base-weight}
\end{equation}

基础等权只是可审计起点，不是最终目标。即使某只证券得分特别高，也不能临时提高其基础权重。本章既不重新优化第 8 章因子权重，也不覆盖 \texttt{combined\_score}；基础和最终权重必须分列保存。

\section{个股最大权重}

多头、不加杠杆和个股上限可写为

\begin{align}
  w_{i,t} &\ge 0, \\
  \sum_i w_{i,t} &\le 1, \\
  w_{i,t} &\le w_{\max}.
  \label{eq:ch09-long-cap}
\end{align}

\(w_{\max}\) 来自冻结配置。若基础等权已经超过上限，说明候选数量和约束不能同时充分投资。截断超限权重后，任何重新分配都必须再次检查上限；不能先截断，再无条件归一化到 1，导致原本合格的证券重新超限。

\section{行业偏离约束}

令 \(g(i,t)\) 为信号日已生效的历史行业，组合行业权重为

\begin{equation}
  W_{g,t}=\sum_{i:g(i,t)=g}w_{i,t}.
  \label{eq:ch09-industry-weight}
\end{equation}

若存在信号日可得的基准行业权重 \(b_{g,t}\)，则可要求

\begin{equation}
  \left|W_{g,t}-b_{g,t}\right|\le \delta_g.
  \label{eq:ch09-industry-deviation}
\end{equation}

历史行业和基准版本都不得晚于信号日。若没有可靠历史基准权重，可采用预注册行业配额或行业内选股作为教学替代，但必须使用独立约束版本。行业约束可能削弱高分证券的表达，这是需要报告的约束代价。

\section{流动性约束}

流动性只能使用信号日前已可得的历史代理。一般目标名义金额上限可写为

\begin{equation}
  \operatorname{TargetNotional}_{i,t}
  \le
  \operatorname{ParticipationCap}
  \times
  \operatorname{HistoricalLiquidity}_{i,t}.
  \label{eq:ch09-liquidity-cap}
\end{equation}

历史窗口必须截止于信号日。资金规模、参与率、价格与成交额口径需实际确认。本章只把代理转换为目标权重上限；下一交易日是否真的成交由第 10 章处理，不能使用下一交易日成交量或不可买状态反向限制目标。

\section{换手预算}

设上一期实际证券权重为 \(h^-_{i,t}\)，目标权重为 \(w_{i,t}\)。教学口径的单边目标换手定义为

\begin{equation}
  \operatorname{Turnover}^{\mathrm{target}}_t
  = \frac{1}{2}\sum_i\left|w_{i,t}-h^-_{i,t}\right|.
  \label{eq:ch09-target-turnover}
\end{equation}

上一期实际持仓在信号日是已知状态，可以作为输入；目标换手不是实际成交换手。若超过冻结预算 \(B_t\)，可按固定比例收缩目标变化或保留部分旧持仓，随后重新检查个股、行业和流动性硬约束。下一交易日未成交信息不能改变本章目标换手。

\section{持仓缓冲与进入/退出阈值}

缓冲区让新证券只有进入更高阈值才买入，让已持有证券只有跌出较低阈值才退出；进入阈值高于退出阈值，可减少边界附近反复交易。缓冲会让目标组合依赖上一期实际持仓，但这不是未来信息。

缓冲参数必须在回测前冻结，并比较加入前后的信号表达、目标换手和集中度。本章将其作为可选教学方案，不作为主基线。不能根据第 10 章的换手和净收益反复调整阈值。

\section{现金权重}

证券目标权重之和不必在不可行时强制等于 1。现金权重定义为

\begin{equation}
  w^{\mathrm{cash}}_t=1-\sum_i w_{i,t},\qquad w^{\mathrm{cash}}_t\ge 0.
  \label{eq:ch09-cash-weight}
\end{equation}

现金是约束后的合法结果，不是自动报错。若个股、行业、流动性或换手限制阻止充分投资，应增加现金并记录 \texttt{PARTIAL\_ALLOCATION}，而不是放宽硬约束或强行归一化。

\section{多项约束的优先级}

\begin{table}[htbp]
  \centering
  \small
  \caption{规则型组合构建的约束优先级}
  \label{tab:ch09-priority}
  \begin{tabularx}{\textwidth}{p{1.4cm} p{4.6cm} Y}
    \toprule
    优先级 & 约束或目标 & 违反时处理 \\
    \midrule
    1 & 数据、主键和版本合法 & 立即失败，不生成目标 \\
    2 & 不做空、不加杠杆 & 硬约束，不放宽 \\
    3 & 个股权重上限 & 缩减并再次验证 \\
    4 & 已确认行业约束 & 按固定顺序补充或缩减 \\
    5 & 已确认流动性上限 & 降低目标，不使用未来成交数据 \\
    6 & 换手预算 & 比例收缩或保留旧持仓 \\
    7 & 尽可能表达综合分排序 & 在硬约束可行域内执行 \\
    8 & 充分投资 & 不能压倒更高优先级约束 \\
    9 & 剩余部分持有现金 & 记录部分配置状态 \\
    \bottomrule
  \end{tabularx}
\end{table}

同时适用的基础约束汇总见 \cref{tab:ch09-constraint-summary}。

\begin{table}[htbp]
  \centering
  \footnotesize
  \caption{个股、行业、流动性和换手约束摘要}
  \label{tab:ch09-constraint-summary}
  \begin{tabularx}{\textwidth}{p{2.5cm} p{4.0cm} p{3.6cm} Y}
    \toprule
    约束 & 信号日输入 & 目标检查 & 禁止使用 \\
    \midrule
    个股 & 冻结 \(w_{\max}\) & 每只证券目标权重 & 历史收益择优上限 \\
    行业 & 点时行业、点时基准或配额 & 行业权重与偏离 & 未来行业或基准版本 \\
    流动性 & 信号日前历史代理与资金规模 & 目标名义金额上限 & 下一日成交量和成交状态 \\
    换手 & 上一期实际持仓、冻结预算 & 目标差额口径 & 实际成交或未成交信息 \\
    \bottomrule
  \end{tabularx}
\end{table}

\section{约束冲突、不可行与确定性降级}

约束冲突不能靠随机删股或静默放宽解决。所有降级动作都要按配置顺序执行，并记录输入状态、触发原因、调整前后权重和仍未满足的约束。

\begin{figure}[htbp]
  \centering
  \begin{tikzpicture}[
    node distance=4mm,
    every node/.style={font=\small},
    quantPortfolioDecision/.style={diamond, aspect=2.4, draw=WarmOrange, fill=WarmOrange!8, align=center, inner sep=2pt},
    quantPortfolioAction/.style={quantCard, text width=7.2cm, minimum height=0.9cm, align=center},
    quantPortfolioSuccess/.style={quantGreenCard, text width=7.2cm, minimum height=0.9cm, align=center},
    quantPortfolioFailure/.style={quantOrangeCard, text width=7.2cm, minimum height=0.9cm, align=center}
  ]
    \node[quantPortfolioDecision] (feasible) {全部硬约束可行？};
    \node[quantPortfolioSuccess, below left=7mm and 20mm of feasible] (success) {是：冻结目标权重与约束报告};
    \node[quantPortfolioAction, below right=7mm and 20mm of feasible] (reduce) {否：按高分顺序补充或缩减超限权重};
    \node[quantPortfolioAction, below=of reduce] (turnover) {按冻结规则收缩换手并保留必要旧持仓};
    \node[quantPortfolioAction, below=of turnover] (cash) {无法充分分配的部分转为现金};
    \node[quantPortfolioDecision, below=of cash] (recheck) {重新验证硬约束？};
    \node[quantPortfolioSuccess, below left=7mm and 20mm of recheck] (partial) {通过：\texttt{PARTIAL\_ALLOCATION}};
    \node[quantPortfolioFailure, below right=7mm and 20mm of recheck] (failed) {仍失败：\texttt{INFEASIBLE}};

    \draw[quantArrow] (feasible) -- node[above left]{是} (success);
    \draw[quantArrow] (feasible) -- node[above right]{否} (reduce);
    \draw[quantArrow] (reduce) -- (turnover);
    \draw[quantArrow] (turnover) -- (cash);
    \draw[quantArrow] (cash) -- (recheck);
    \draw[quantArrow] (recheck) -- node[above left]{是} (partial);
    \draw[quantArrow] (recheck) -- node[above right]{否} (failed);
  \end{tikzpicture}
  \caption{约束不可行时的确定性降级决策树（结构示意）}
  \label{fig:ch09-degradation-tree}
\end{figure}

允许的动作包括从下一名高分证券补充、缩减超限权重、按固定顺序重新分配、比例收缩交易、保留部分旧持仓和增加现金。任何动作后都必须重复验证高优先级硬约束。

\section{目标权重与期望交易量}

目标交易差额定义为

\begin{equation}
  \operatorname{desired\_trade\_weight}_{i,t}
  = w_{i,t}-h^-_{i,t}.
  \label{eq:ch09-desired-trade}
\end{equation}

正值表示期望增加，负值表示期望减少，零表示目标不变。若需要目标名义金额，可在冻结资金规模下转换，但仍不能称为已下单或已成交。本章不生成 \texttt{filled\_weight} 或 \texttt{actual\_holding\_after\_trade}。

\section{组合约束检查与审计表}

\begin{table}[htbp]
  \centering
  \small
  \caption{目标组合约束检查模板}
  \label{tab:ch09-constraint-audit}
  \begin{tabularx}{\textwidth}{p{3.2cm} p{3.1cm} p{3.1cm} Y}
    \toprule
    检查项 & 冻结限制 & 目标观察值 & 状态与降级记录 \\
    \midrule
    证券权重非负 & \(w_{i,t}\ge0\) & -- & -- \\
    证券权重与现金 & 合计为 1 & -- & -- \\
    最大个股权重 & \(w_{\max}\) & -- & -- \\
    行业偏离 & \(\delta_g\) 或配额 & -- & -- \\
    流动性上限 & 历史代理映射上限 & -- & -- \\
    目标换手 & \(B_t\) & -- & -- \\
    证券数量与覆盖 & 配置与状态规则 & -- & -- \\
    组合版本 & 唯一且不可覆盖 & -- & -- \\
    \bottomrule
  \end{tabularx}
\end{table}

审计报告还应保存各行业权重及偏离、流动性覆盖、约束失败、每步重新分配、现金来源和最终 \texttt{portfolio\_status}。只保留最终权重会丢失约束如何影响信号的证据链。

\section{一个匿名教学案例}

\begin{casebox}
\textbf{仅用于解释选股、权重和约束逻辑，不代表真实证券、真实数据库、真实投资组合或实证研究结果。}

匿名信号日 \(t^\star\) 有九只证券，教学配置固定选择前四名。上一期证券权重合计为 0.80，现金为 0.20。案例使用匿名约束：个股上限、行业配额和流动性目标上限只为解释计算，不代表真实团队阈值。

\begin{center}
\footnotesize
\begin{tabular}{c c r r c r}
  \toprule
  证券 & 行业 & 综合分 & 上期实际权重 & 前四名 & 约束提示 \\
  \midrule
  A & 甲 & 2.1 & 0.15 & 是 & 高分，受个股上限 \\
  B & 甲 & 1.8 & 0.10 & 是 & 行业甲接近上限 \\
  C & 乙 & 1.6 & 0.08 & 是 & 流动性目标上限较低 \\
  D & 丙 & 1.2 & 0.12 & 是 & 个股上限 \\
  E & 乙 & 0.9 & 0.10 & 否 & 退出目标 \\
  F & 丙 & 0.7 & 0.08 & 否 & 退出目标 \\
  G & 甲 & 0.4 & 0.07 & 否 & 退出目标 \\
  H & 乙 & 0.1 & 0.05 & 否 & 退出目标 \\
  I & 丙 & -0.2 & 0.05 & 否 & 退出目标 \\
  \bottomrule
\end{tabular}
\end{center}

前四名的基础等权均为 0.25。确定性约束后，A 因个股上限降至 0.20；行业甲总额限制使 B 只能为 0.10；C 的历史流动性映射上限为 0.10；D 受个股上限约束为 0.20。证券目标权重合计 0.60，现金为 0.40。

\begin{center}
\footnotesize
\begin{tabular}{c r r r r}
  \toprule
  证券 & 基础权重 & 上期实际权重 & 目标权重 & 期望交易差额 \\
  \midrule
  A & 0.25 & 0.15 & 0.20 & 0.05 \\
  B & 0.25 & 0.10 & 0.10 & 0.00 \\
  C & 0.25 & 0.08 & 0.10 & 0.02 \\
  D & 0.25 & 0.12 & 0.20 & 0.08 \\
  E & 0.00 & 0.10 & 0.00 & -0.10 \\
  F & 0.00 & 0.08 & 0.00 & -0.08 \\
  G & 0.00 & 0.07 & 0.00 & -0.07 \\
  H & 0.00 & 0.05 & 0.00 & -0.05 \\
  I & 0.00 & 0.05 & 0.00 & -0.05 \\
  \bottomrule
\end{tabular}
\end{center}
\end{casebox}

案例的目标单边换手按 \cref{eq:ch09-target-turnover} 为 0.25。表中目标权重和期望差额均不代表下一交易日实际结果；若某笔交易无法成交，第 10 章会保留旧实际持仓并记录偏离，而不是回写本表。

\section{pandas 风格最小实现}

\cref{lst:ch09-portfolio-builder} 展示规则型目标组合构建的职责边界。代码只生成目标和审计状态，不生成订单成交或回测结果。

\begin{lstlisting}[
  caption={约束型多头目标组合的最小接口},
  label={lst:ch09-portfolio-builder}
]
FORBIDDEN_COLUMNS = {
    "execution_status", "filled", "execution_price",
    "limit_status", "next_day_suspension",
    "future_return", "backtest_result", "orders", "fills",
}

def build_constrained_long_only_target(
    frozen_scores,
    signal_universe,
    previous_actual_holdings,
    portfolio_config,
):
    assert portfolio_config.combination_version_frozen
    assert (
        frozen_scores["combination_version"].unique().tolist()
        == [portfolio_config.combination_version]
    )
    for frame in (
        frozen_scores,
        signal_universe,
        previous_actual_holdings,
    ):
        assert FORBIDDEN_COLUMNS.isdisjoint(frame.columns)

    assert_unique(frozen_scores, [
        "security_id", "signal_date", "combination_version"
    ])
    validate_point_in_time_columns(
        signal_universe,
        asof_columns=[
            "industry_asof", "benchmark_weight_asof",
            "liquidity_window_end",
        ],
        not_after="signal_date",
    )

    conditions_met = evaluate_frozen_applicability_conditions(
        frozen_scores["applicability_conditions"],
        signal_universe,
        asof="signal_date",
    )
    eligible = frozen_scores.loc[
        frozen_scores["joint_valid"] & conditions_met
    ].copy()
    ranked = deterministic_rank(
        eligible,
        by=["combined_score", "security_id"],
        ascending=[False, True],
        tie_rule=portfolio_config.tie_rule,
    )
    selected = select_from_frozen_rule(
        ranked,
        method=portfolio_config.selection_method,
        count=portfolio_config.fixed_count,
    )
    if selected.empty:
        return append_infeasible_portfolio(
            reason="NO_ELIGIBLE_CANDIDATE",
            portfolio_version=hash_frozen_config(portfolio_config),
        )
    ranked["selected"] = ranked["security_id"].isin(
        selected["security_id"]
    )
    ranked["base_weight"] = 0.0
    ranked.loc[ranked["selected"], "base_weight"] = (
        1.0 / len(selected)
    )

    state = build_portfolio_state_union(
        ranked,
        signal_universe,
        previous_actual_holdings,
        stable_key="security_id",
        include_previous_only_holdings=True,
    )
    state = apply_optional_buffer(
        state,
        enabled=portfolio_config.buffer_enabled,
        entry_rule=portfolio_config.entry_rule,
        exit_rule=portfolio_config.exit_rule,
    )
    state = apply_individual_caps(
        state, cap=portfolio_config.individual_cap
    )
    state = apply_industry_constraints(
        state,
        limits=portfolio_config.industry_limits,
        deterministic_order=ranked["security_id"],
    )
    state = apply_liquidity_caps(
        state,
        capital=portfolio_config.capital,
        participation_cap=portfolio_config.participation_cap,
    )
    state = redistribute_with_frozen_order(
        state,
        ranked=ranked,
        allow_cash=True,
    )
    validate_hard_constraints(state)

    state = shrink_to_turnover_budget(
        state,
        previous_weight="previous_actual_weight",
        budget=portfolio_config.turnover_budget,
        trace=True,
    )
    state = deterministic_degradation(
        state,
        priority=portfolio_config.constraint_priority,
        allow_cash=True,
        statuses=["PARTIAL_ALLOCATION", "INFEASIBLE"],
    )
    validate_hard_constraints(state)

    state["target_weight"] = state["constrained_weight"]
    state["desired_trade_weight"] = (
        state["target_weight"]
        - state["previous_actual_weight"]
    )
    cash_weight = 1.0 - state["target_weight"].sum()
    assert (state["target_weight"] >= 0).all()
    assert cash_weight >= 0
    assert abs(state["target_weight"].sum() + cash_weight - 1) < 1e-12
    assert FORBIDDEN_COLUMNS.isdisjoint(state.columns)

    output = attach_portfolio_audit_fields(
        state,
        cash_weight=cash_weight,
        portfolio_version=hash_frozen_config(portfolio_config),
    )
    portfolio_store.append_only(output)
    return output
\end{lstlisting}

真实实现还应保证候选不足时不发生除零，并让所有降级步骤追加到审计表。复杂二次规划不是本章要求；规则型流程的价值在于每次调整都能重放和解释。

\section{自动测试与人工核对}

自动测试至少覆盖：

\begin{enumerate}
  \item 非冻结 \texttt{combination\_version} 不能使用；
  \item \texttt{joint\_valid = False} 的记录不能进入候选；
  \item 相同证券、信号日和版本重复时失败；
  \item 排序不能读取未来标签；
  \item 并列处理确定且可重复；
  \item 固定数量或固定分位规则严格来自配置；
  \item 任一证券目标权重不得为负；
  \item 证券目标权重与现金权重合计为 1；
  \item 任一个股不得超过上限；
  \item 重新分配后再次检查个股上限；
  \item 历史行业区间覆盖信号日；
  \item 基准行业权重版本不晚于信号日；
  \item 流动性窗口不包含未来数据；
  \item 下一交易日成交量或状态不得参与目标构建；
  \item 目标换手与实际成交换手使用不同字段名；
  \item 换手预算降级过程可追溯；
  \item 无法充分投资时现金权重为正；
  \item 不可行状态不能被静默忽略；
  \item \texttt{target\_weight}、\texttt{desired\_trade\_weight} 与 \texttt{previous\_actual\_weight} 可相互还原；
  \item 组合版本变化保留旧输出；
  \item 相同输入、配置和版本得到相同结果；
  \item 目标组合函数不能访问 \texttt{orders}、\texttt{fills}、\texttt{execution\_price} 或 \texttt{backtest\_result}。
\end{enumerate}

\begin{practicebox}
\textbf{[正确实践]} 人工至少核对一个普通候选、一个并列边界、一个个股超限、一个行业受限、一个流动性受限、一个换手收缩和一个现金配置案例。逐项重算基础权重、每步约束、目标权重、期望差额和组合状态，并确认没有读取下一交易日字段。
\end{practicebox}

\section{常见错误}

\begin{warningbox}
常见错误包括：把综合分直接当权重；根据第 10 章收益选择持股数量；使用下一交易日不可买状态重排信号日候选；用未来基准权重或行业分类；先截断个股权重再归一化导致重新超限；行业约束后不复查个股约束；为满仓违反流动性限制；把目标换手称为实际换手；把目标组合称为实际持仓；无法卖出时让旧持仓凭空消失；不可行时静默放宽约束；只保存最终权重而不保存基础权重和降级过程；把现金视为回测错误；以及使用组合结果回头修改第 8 章因子集合或权重。
\end{warningbox}

这些错误会把信号、决策和执行三层混在一起，使目标组合无法审计。正确修复是恢复时间边界和状态账本，不是借未来执行信息重写信号日决策。

\section{本章产出}

\begin{outputbox}
本章应提交：冻结候选规则与并列规则；基础等权；个股、行业、流动性和换手约束配置；缓冲方案状态；约束优先级和确定性降级记录；证券目标权重、期望交易差额与现金权重；组合约束检查；自动测试与人工核对结果；以及进入第 10 章的唯一冻结 \texttt{portfolio\_version}。
\end{outputbox}

\begingroup
\small
\begin{longtable}{@{}p{4.4cm}p{4.8cm}p{5.6cm}@{}}
\caption{目标组合证券级输出 schema}\label{tab:ch09-output-schema}\\
\toprule
字段 & 含义 & 审计要求 \\
\midrule
\endfirsthead
\toprule
字段 & 含义 & 审计要求 \\
\midrule
\endhead
\texttt{security\_id}、\texttt{signal\_date} & 稳定主键与信号日 & 同一组合版本唯一 \\
\texttt{combination\_version} & 第 8 章冻结综合分版本 & 不得重选历史版本 \\
\texttt{portfolio\_version} & 目标组合版本 & 规则或约束变化即新版本 \\
\texttt{combined\_score}、\texttt{score\_rank} & 冻结分数与确定性排名 & 权重不能覆盖分数 \\
\texttt{selected}、\texttt{selection\_reason} & 入选状态与原因 & 未入选记录仍保留 \\
\texttt{base\_weight} & 约束前等权起点 & 与候选数量一致 \\
\texttt{previous\_actual\_weight} & 上一期实际证券权重 & 来自信号日已知快照 \\
\texttt{target\_weight} & 约束后目标权重 & 非负且满足硬约束 \\
\texttt{desired\_trade\_weight} & 目标减上期实际权重 & 不代表成交 \\
\texttt{industry\_asof} & 点时行业 & 生效日不晚于信号日 \\
\texttt{benchmark\_industry\_weight\_asof} & 点时基准行业权重 & 无可靠版本时使用登记替代 \\
\texttt{liquidity\_proxy\_asof} & 历史流动性代理 & 窗口截止信号日 \\
\texttt{individual\_cap} & 个股上限 & 来自冻结配置 \\
\texttt{industry\_constraint\_status} & 行业约束状态 & 保存偏离与调整 \\
\texttt{liquidity\_constraint\_status} & 流动性约束状态 & 不使用未来成交量 \\
\texttt{turnover\_constraint\_status} & 换手预算状态 & 保存收缩轨迹 \\
\texttt{portfolio\_status} & 成功、部分配置或不可行 & 不静默忽略失败 \\
\texttt{cash\_weight} & 组合级现金权重 & 与证券权重合计为 1 \\
\texttt{constraint\_version} & 约束配置版本 & 可重放优先级与降级 \\
\bottomrule
\end{longtable}
\endgroup

组合级汇总另保存证券目标权重总和、现金、目标换手、行业权重及偏离、最大个股权重、证券数量、流动性覆盖和所有失败/降级记录。

\section{章节自测}

\begin{quizbox}
\begin{enumerate}
  \item 综合得分、目标组合、期望订单、实际成交和实际持仓分别表示什么？
  \item 为什么本章选择固定数量作为主基线，而固定分位只作敏感性方案？
  \item 基础等权为什么不能直接称为最终目标权重？
  \item 截断个股权重后为什么不能无条件归一化到满仓？
  \item 行业与流动性数据的点时边界是什么？
  \item 目标换手与实际成交换手为何必须分列？
  \item 持仓缓冲如何降低交易，又引入了什么路径依赖？
  \item 哪些情况下现金权重为正是合法结果？
  \item 不可行时，确定性降级为何优于随机删股或静默放宽？
  \item 进入第 10 章前必须冻结哪些版本、状态和审计记录？
\end{enumerate}
\end{quizbox}

\section{与第 10 章《含成本回测与结果解释》的衔接}

本章交付的是冻结目标组合与期望交易差额，不是下一交易日实际持仓。第 10 章将按事件顺序生成订单、判断停牌和价格限制、模拟成交、计算费用、更新实际持仓并解释目标偏离；它必须保留无法成交的旧持仓，不能回写本章候选或目标权重。

交接包至少包含 \texttt{portfolio\_version}、证券目标权重、期望差额、现金、上一期实际持仓快照、约束状态、降级轨迹和点时数据版本。只有这些输入全部冻结后，才能开始含成本回测。
